\subsection*{\texttt{chunkat(...)}：等分}

\begin{lstlisting}[language=croqtile,caption={chunkat(...)：等分}]
lhs.chunkat(tr, tc)
\end{lstlisting}

表达式 \co@chunkat@ 根据对应索引的上界划分各维。输入为 \co@[64, 128]@、并行声明为 \co@parallel {tr, tc} by [4, 8]@ 时，选中的分块大小是 \co@[16, 16]@。\co@tr@ 和 \co@tc@ 的当前值决定选哪块，上界决定块有多大。

已有并行或循环索引明确了分块方式时，可以用 \co@chunkat@ 选取当前分块。本例中，\co@lhs@ 的形状为 \co@[64, 128]@，\co@tr@、\co@tc@ 的范围分别为 \co@4@、\co@8@，所以每块为 \co@[16, 16]@。索引当前值再决定选中哪一块。

\begin{figure}[htbp]\centering
\includegraphics[width=\linewidth,height=0.56\textheight,keepaspectratio]{f04-fig3_chunkat_light.png}
\caption{对按 [4, 8] 网格划分的 [64, 128] 张量，chunkat(1, 3) 选中第 16–31 行、第 48–63 列。}
\end{figure}
